% ANCHOR Dependencies
\RequirePackage{listings}
\RequirePackage[ruled,vlined]{algorithm2e}

% ANCHOR Mono Font for Code
\newfontfamily{\FiraCode}{Fira Code}
\newfontfamily{\IBMMono}{IBM Plex Mono}
\newfontfamily{\OperatorMono}{Operator Mono Book}
\newfontfamily{\SFMono}{SF Mono Regular}
\newfontfamily{\Consolas}{Consolas}
\newfontfamily{\TimesRoman}{Cambria}

\lstdefinestyle{Fira}{basicstyle=\FiraCode}
\lstdefinestyle{IBM}{basicstyle=\IBMMono}
\lstdefinestyle{Operator}{basicstyle=\OperatorMono}
\lstdefinestyle{SF}{basicstyle=\SFMono}
\lstdefinestyle{Basic}{basicstyle=\large\Consolas}

\newcommand{\codeit}[1]{
   \texorpdfstring{\textit{\OperatorMono #1}}{#1}
}

% ANCHOR C++ code listing style
\definecolor{kword}{HTML}{db2777}
\lstdefinestyle{cpp}{
   language=C++,
   basicstyle=\linespread{0.9}\ttfamily,
   keywordstyle = {\color{kword}},
   keywordstyle = \texttt,
   otherkeywords = {<,>,+,-,&,|,!,=,?, :},
   breaklines=true,
   tabsize=3,
   frame=lines,
   rulecolor=\color{gray!50},
   % backgroundcolor = \color{lightgray!10}
}
\lstnewenvironment{cpp}[1][]{\lstset{style=cpp}}{}

% ANCHOR Pseudo code listing style
\lstdefinelanguage{pseudo}{
  morekeywords={BEGIN,END,IF,ELSE,LOOP,FOREVER,READ,DISPLAY,WAIT,TURN,ON,OFF,CLEAR,LOG,DELAY},
  sensitive=false,
  morecomment=[l]{//},
  morecomment=[l]{===},
}
\lstdefinestyle{pseudo}{
  language=pseudo,
  basicstyle=\ttfamily\small,
  keywordstyle=\ttfamily,
  commentstyle=\color{comment}\vspace{0.025em},
  showstringspaces=false,
  breaklines=true,
  frame=none,
  columns=fullflexible
}
\lstnewenvironment{pseudocode}[1][]{\lstset{style=pseudo}}{}


% ANCHOR Arduino code listing style
\lstdefinelanguage{arduino}{
  language=C++,
  morekeywords={
   oled,
  },
  morekeywords=[2]{
    pinMode, digitalWrite, digitalRead, analogRead, analogWrite,
    delay, millis, attachInterrupt, detachInterrupt,
    setup, loop, clearDisplay, setCursor, print, println, readTemperature, readHumidity, readPH
  },
  morekeywords=[3]{INPUT, OUTPUT, HIGH, LOW, NULL, millis, delayMicroseconds},
  sensitive=true,
  morecomment=[l]{//},
  morecomment=[s]{/*}{*/},
  morestring=[b]",
}
\lstdefinestyle{arduinoStyle}{
   language=arduino,
   basicstyle=\ttfamily\small\setstretch{1.3},
   keywordstyle=\color{keyword},
   keywordstyle=[2]\color{function},
   keywordstyle=[3]\color{number},
   commentstyle=\color{comment}\itshape,
   stringstyle=\color{string},
   showstringspaces=false,
   breaklines=true,
   frame=none,
   columns=fullflexible,
}
\lstnewenvironment{arduino}[1][]{\lstset{style=arduinoStyle}}{}

\newcommand{\includecode}[1]{\lstinputlisting[style=cpp]{#1}}
\newcommand{\code}[1]{\texttt{#1}}
\newtcbox{\codebox}[1][gray!18]{%
    on line, 
    box align=base, 
    nobeforeafter, 
    colback=#1,  % Background color
    colframe=white, % Set the border color to white
    opacityframe=0, % Make the border completely transparent
    boxrule=0pt, 
    arc=2pt, 
    left=2pt, 
    right=2pt, 
    top=3pt, 
    bottom=3pt, 
    boxsep=0pt
}
\newcommand{\bcode}[2][gray!15]{\codebox[#1]{\texttt{#2}}}